% Werke X/1 (1917), native152–155, printed145–148.
\begingroup
\setlength{\headheight}{25pt}
\SourceLeaf{152}\NoSourceHead
\HistoricTitle{[ÄLTESTE UNTERSUCHUNGEN\\ ÜBER LEMNISKATISCHE FUNKTIONEN.]}
\SourceGroup{I}
\begin{center}\bfseries [REIHENENTWICKELUNGEN UND ADDITIONSTHEOREME.]\end{center}
\renewcommand{\thefootnote}{*)}
\HistoricCitation{[Eintragungen im \textsc{Leiste}\footnotemark.]}
\renewcommand{\thefootnote}{[*)}\footnotetext{Vgl. die Fußnote auf S. 78.]}
\NumberedItem{1}
\noindent{\small [Einbanddecke]}
\[
\int\frac{dp}{\surd(1-p^4)}=\left[\frac{\pi}{4}\right]\frac{2}{\int\surd\sin x\,dx}=1{,}311031
\]
\noindent nach \textsc{Stirling}: \textit{De summatione et interpolatione serierum}
\renewcommand{\thefootnote}{[**)]}\footnotemark
\renewcommand{\thefootnote}{[**)}\footnotetext{\textit{Methodus differentialis sive Tractatus de summatione et interpolatione serierum infinitarum}, auctore \textsc{Jacobo Stirling}, Londini 1730, S. 58.]}
\[\begin{gathered}1{,}31102877714605987\\[6pt]1{,}2453.\end{gathered}\]
\ClosingRule
\vfill\noindent X1.\hfill19\par
\clearpage
\SourceLeaf{153}\ArithmeticHead{146}{L}{ANALYSIS.\quad NACHLASS.}
\NumberedItem{2}
\noindent{\small [Einbanddecke]}
\[\setlength{\arraycolsep}{3pt}\renewcommand{\arraystretch}{1.65}
\begin{array}{rrrrrl}
1[\varphi]&-\dfrac1{10}[\varphi^5]&-\dfrac1{24}[\varphi^9]&-\dfrac5{208}[\varphi^{13}]&-\dfrac{35}{17.128}[\varphi^{17}]&\ldots\\
&&+\dfrac1{20}&+\dfrac7{120}&+\dfrac9{208}&\ldots\\
&&&-\dfrac7{200}&+\dfrac1{64}&\ldots\\
&&&&-\dfrac{57}{800}&\ldots\\
&&&&+\dfrac{57}{2000}&\ldots\\\hline
1[\varphi]&-\dfrac1{10}[\varphi^5]&+\dfrac1{120}[\varphi^9]&-\dfrac{11}{15600}[\varphi^{13}]&+\dfrac{211}{3536000}[\varphi^{17}]&\ldots
\end{array}\]
\[\begin{gathered}
\left[\frac1{10}=\right]\frac12\mathpunct{.}\frac15,\qquad
\left[\frac1{120}=\right]\frac38\mathpunct{.}\frac15\mathpunct{.}\frac19,\\[5pt]
\left[\frac{11}{15600}=\right]\frac{33}{80}\frac{1.1.1}{5.9.13},\qquad
\left[\frac{211}{3536000}=\right]\frac{211.9}{3200}\mathpunct{.}\frac{1.1.1.1}{5.9.13.17}.
\end{gathered}\]
\ClosingRule
\NumberedItem{3}
\noindent{\small [Schutzblatt]}
\[\begin{aligned}
\int\surd\sin x\,dx&=\tfrac23x^{\frac23}-\int dx(\surd x-\surd\sin x)\\
&=2\int\frac{pp\,dp}{\surd(1-p^4)}
\end{aligned}\]
\begin{flushright}
$\begin{array}{r}
0{,}7849171\\1363706\\1{,}6337417\\1794738\\1{,}5163437\\2928571\\\hline
4{,}543704
\end{array}\quad\begin{array}{l}\vphantom{0{,}7849171}\\\vphantom{1363706}\\\vphantom{1{,}6337417}\\\vphantom{1794738}\\\vphantom{1{,}5163437}\\\vphantom{2928571}\\\times\dfrac{\pi}{12}=1{,}189539\end{array}$
\end{flushright}
\[\text{\textsc{Fuss}}\quad1{,}198\pm=\frac{\pi}{\varpi}\]
\begin{center}\textsc{Euler} applicata mea reductione $1{,}198122.$\end{center}
\ClosingRule
\clearpage
\SourceLeaf{154}\ArithmeticHead{147}{R}{\scriptsize\shortstack{ÄLTESTE UNTERSUCHUNGEN ÜBER\\LEMNISKATISCHE FUNKTIONEN.}}
\NumberedItem{4}
\noindent{\small [S. 16]}
\par\noindent\hspace*{\parindent}In der Curva (Elastica) Lemniscata ist
\[\begin{array}{ll}s&\text{Sehne eines Bogens}\\c&\text{Sehne des Compl[ements] des B[ogens]}\end{array}\]
\[ss+cc+ccss=1\]
\noindent Sehne des dopp[elten Bogens]
\[\frac{2(s+s^3c}{1+s^4}=\frac{2sc}{1-sscc}=\frac{2sc}{ss+cc}\]
\noindent Cosehne des dopp[elten Bogens]
\[\frac{1-2ss-s^4}{1+2ss-s^4}=\frac{cc-ss}{1+ccss}=\frac{cc-ss}{2-ss-cc}\]
\[\renewcommand{\arraystretch}{1.8}\begin{array}{llll}
\text{Sehne der Summe}&\dfrac{s'c+sc'}{1-s'sc'c}&\text{d[er] Diff[erenz]}&\dfrac{s'c-sc'}{1+s'sc'c}\\
\text{Cosehne der Summe}&\dfrac{c'c-s's}{1+s'sc'c}&\text{d[er] Diff[erenz]}&\dfrac{c'c+s's}{1-s'sc'c}
\end{array}\]
\noindent{\small [S. 17--18 nach längerer Rechnung]}
\par\medskip\noindent[Sehne des doppelten Bogens]
\[sc\frac{2+2ss}{1+s^4}=sc\frac{2+2cc}{1+c^4}\]
\noindent[Cosehne des doppelten Bogens]
\[\frac{1-2ss-s^4}{1+2ss-s^4}=\frac{1-2cc-c^4}{1+2cc-c^4}\]
\noindent[Sehne des dreifachen Bogens]
\[s\frac{3-6s^4-s^8}{1+6s^4-3s^8}\]
\noindent[Cosehne des dreifachen Bogens]
\[c\frac{3-6c^4-c^8}{1+4cc-6c^4+4c^6+c^8}=c\frac{1+6s^4-3s^8}{1+4ss-6s^4+4s^6+s^8}.\]
\ClosingRule
\clearpage
\NumberedItem{5}
\noindent{\small [S. 20]}
\par\noindent\hspace*{\parindent}Theilt man also auch hier die ganze Peripherie in $360^0$, so ist
\[\setlength{\arraycolsep}{5pt}\renewcommand{\arraystretch}{1.65}\begin{array}{crl}
\text{chorda}&&\\\text{arcus}&&\\
0&0&\\
30^0&&=\mathop{\surd}\nolimits\dfrac{1-\surd(\surd12-3)}{1+\surd(\surd12-3)}=\dfrac{1-\surd(\surd12-3)}{\surd(4-\surd12)}\\
45^0&0{,}6435943&=\surd(\surd2-1)\\
60^0&0{,}8253788&={}^4\!\surd(\surd12-3)=\surd(\surd\tfrac32-\surd\tfrac12){}^8\!\surd3\\
90^0&1&
\end{array}\]
\ClosingRule
\vfill\noindent\hfill19\textsuperscript{*}\par
\clearpage
\SourceLeaf{155}\ArithmeticHead{148}{L}{ANALYSIS.\quad NACHLASS.}
\NumberedItem{6}
\noindent{\small [S. 20--21]}
\[\begin{aligned}
d\sin&=\frac{2}{c+\frac1c}d\varphi=\frac{1+ss}{cs}d\varphi\\[5pt]
d\cos&=\frac{-2}{s+\frac1s}d\varphi\\[5pt]
\sin&=\varphi-\frac1{10}\varphi^5+\frac1{120}\varphi^9+\ldots\\
\cos&=1-\varphi\varphi+\tfrac12\varphi^4-\ldots
\end{aligned}\]
\[\begin{aligned}
\frac1{\sin}&=\frac1\varphi+\frac1{10}\varphi^3+\ldots\\
\frac1{\cos}&=1+\varphi\varphi+\frac12\varphi^4+\ldots
\end{aligned}\]
\[\varphi=s+\frac12\mathpunct{.}\frac15s^5+\frac12\mathpunct{.}\frac34\mathpunct{.}\frac19s^9+\frac{1.3.5}{2.4.6}\mathpunct{.}\frac1{13}s^{13}+\ldots\]
\noindent Man nenne nach der Analogie des Kreises
\[\frac sc\ldots t,\]
\noindent so ist
\[dt=\left(\frac12\left(1+\frac1{cc}\right)+\frac12\left(\frac1{cc}+\frac{ss}{cc}\right)\right)d\varphi.\]
\[s\,4=sc\frac{1+ss}{1+s^4}\frac{[4(1-5s^4-5s^8+s^{12})(1+s^4)]}{1+20s^4-26s^8-20s^{12}+s^{16}}.\]
\ClosingRule
\NumberedItem{7}
\noindent{\small S. 24]}
\par\noindent\hspace*{\parindent}Ad p. 16
\[\begin{aligned}
\varphi1&=\varphi\surd\tfrac12+\varphi\surd\tfrac13\\[4pt]
\varphi\surd\tfrac12-\varphi\surd\tfrac13&=\varphi\tfrac17\\[4pt]
\varphi\surd\tfrac12-\varphi\surd\tfrac13&=\varphi1-\varphi\surd\tfrac22\tfrac45\\[4pt]
\varphi1&=\varphi\surd\frac1m+\varphi\surd\frac{m-1}{m+1}\\[4pt]
\varphi\surd\frac1m-\varphi\surd\frac{m-1}{m+1}&=\varphi\frac{mm-2m-1}{mm+2m-1}.
\end{aligned}\]
\ClosingRule
\clearpage
\endgroup
